\ifdefined\gkver\else\def\gkver{student}\fi
\documentclass[\gkver]{gaokaozhenti}
\begin{document}

\chapter{2003年上海理综}
%% paperType: 理综物理部分

\begin{enumerate}
\item
%% number: 10
%% paperName: 2003年上海理综第 10 题 3 分
%% typeId: 1
%% type: 单选题
%% chapter: 电路
%% point: 传感器
%% method: 无
%% score: 3
%% degree: 650
%% duplicateId: 0
%% body:
演示位移传感器的工作原理如右图示，物体 $M$ 在导轨上平移时，带动滑动变阻器的金属滑杆 $P$，通过电压表显示的数据，来反映物体位移的大小 $x$。假设电压表是理想的，则下列说法正确的是\xzanswer{B}

\onepicture[w=5.4cm]{figs/10.png}

\fourchoices[answer=B]
{物体 $M$ 运动时，电源内的电流会发生变化}
{物体 $M$ 运动时，电压表的示数会发生变化}
{物体 $M$ 不动时，电路中没有电流}
{物体 $M$ 不动时，电压表没有示数}

%% answer: B
%% memo:
\memoanswer{
本题原卷及材料中未提供详解，待补充。
}

\item
%% number: 11
%% paperName: 2003年上海理综第 11 题 3 分
%% typeId: 1
%% type: 单选题
%% chapter: 电磁感应
%% point: 电磁感应现象
%% method: 无
%% score: 3
%% degree: 550
%% duplicateId: 0
%% body:
唱卡拉 OK 用的话筒，内有传感器。其中有一种是动圈式的，它的工作原理是在弹性膜片后面粘接一个轻小的金属线圈，线圈处于永磁体的磁场中，当声波使膜片前后振动时，就将声音信号转变为电信号。下列说法正确的是\xzanswer{B}

\fourchoices[answer=B]
{该传感器是根据电流的磁效应工作的}
{该传感器是根据电磁感应原理工作的}
{膜片振动时，穿过金属线圈的磁通量不变}
{膜片振动时，金属线圈中不会产生感应电动势}

%% answer: B
%% memo:
\memoanswer{
本题原卷及材料中未提供详解，待补充。
}

\item
%% number: 12
%% paperName: 2003年上海理综第 12 题 3 分
%% typeId: 1
%% type: 单选题
%% chapter: 电路
%% point: 传感器
%% method: 无
%% score: 3
%% degree: 850
%% duplicateId: 0
%% body:
用遥控器调换电视机频道的过程，实际上就是传感器把光信号转化为电信号的过程。下列属于这类传感器的是\xzanswer{A}

\fourchoices[answer=A]
{红外报警装置}
{走廊照明灯的声控开关}
{自动洗衣机中的压力传感装置}
{电饭煲中控制加热和保温的温控器}

%% answer: A
%% memo:
\memoanswer{
本题原卷及材料中未提供详解，待补充。
}

\item
%% number: 13
%% paperName: 2003年上海理综第 13 题 3 分
%% typeId: 1
%% type: 单选题
%% chapter: 气体性质
%% point: 能源的开发与利用
%% method: 无
%% score: 3
%% degree: 750
%% duplicateId: 0
%% body:
自行车、电动自行车、普通汽车消耗能量的类型分别是\xzanswer{B}

①生物能　②核能　③电能　④太阳能　⑤化学能

\fourchoices[answer=B]
{①④⑤}
{①③⑤}
{①②③}
{①③④}

%% answer: B
%% memo:
\memoanswer{
本题原卷及材料中未提供详解，待补充。
}

\item
%% number: 14
%% paperName: 2003年上海理综第 14 题 3 分
%% typeId: 1
%% type: 单选题
%% chapter: 曲线运动
%% point: 线速度角速度
%% method: 无
%% score: 3
%% degree: 650
%% duplicateId: 0
%% body:
某品牌电动自行车的铭牌如下：

\begin{center}
\begin{tblr}{|l|l|}
\hline
车型：20 吋（车轮直径：$508\Umm$） & 电池规格：$36\UV$、$12\UAh$（蓄电池） \\
\hline
整车质量：$40\Ukg$ & 额定转速：$210\Urmin$（转/分） \\
\hline
外形尺寸：$L\,1800\Umm\times W\,650\Umm\times H\,1100\Umm$ & 充电时间：$2\sim8\Uh$ \\
\hline
电机：后轮驱动、直流永磁式电机 & 额定工作电压/电流：$36\UV/5\UA$ \\
\hline
\end{tblr}
\end{center}

根据此铭牌中的有关数据，可知该车的额定时速约为\xzanswer{C}

\fourchoices[answer=C]
{$15\Ukmh$}
{$18\Ukmh$}
{$20\Ukmh$}
{$25\Ukmh$}

%% answer: C
%% memo:
\memoanswer{
本题原卷及材料中未提供详解，待补充。
}

\item
%% number: 15
%% paperName: 2003年上海理综第 15 题 3 分
%% typeId: 1
%% type: 单选题
%% chapter: 功和能
%% point: 功
%% method: 无
%% score: 3
%% degree: 550
%% duplicateId: 0
%% body:
在交通运输中，常用“客运效率”来反映交通工具的某项效能，“客运效率”表示每消耗单位能量对应的载客数和运送路程的乘积，即客运效率$=\frac{\text{人数}\times\text{路程}}{\text{消耗能量}}$。一个人骑电动自行车，消耗 $1\UMJ$（$10^{6}\UJ$）的能量可行驶 $30\Ukm$，一辆载有 4 人的普通轿车，消耗 $320\UMJ$ 的能量可行驶 $100\Ukm$，则电动自行车与这辆轿车的客运效率之比是\xzanswer{C}

\fourchoices[answer=C]
{$6:1$}
{$12:5$}
{$24:1$}
{$48:7$}

%% answer: C
%% memo:
\memoanswer{
本题原卷及材料中未提供详解，待补充。
}

\item
%% number: 31
%% paperName: 2003年上海理综第 31 题 2 分
%% typeId: 3
%% type: 填空题
%% chapter: 直线运动
%% point: 平均速度和瞬时速度
%% method: 无
%% score: 2
%% degree: 850
%% duplicateId: 0
%% body:
世博会参观者预计有 7000 万人次，交通网络的建设成为关键。目前上海是快的陆上交通工具是连接浦东国际机场和龙阳路地铁站的磁浮列车，它的时速可达 $432\Ukmh$，能在 $7\Umin$ 内行驶 $31\Ukm$ 的全程。该车的平均速率为\tkanswer[2.6cm]{$2.7\times10^{2}$ 或 266}$\Ukmh$。

%% answer: 
%% memo:
\memoanswer{
本题原卷及材料中未提供详解，待补充。
}

\item
%% number: 45
%% paperName: 2003年上海理综第 45 题 5 分
%% typeId: 3
%% type: 填空题
%% chapter: 光学
%% point: 光的电磁说
%% method: 无
%% score: 5
%% degree: 750
%% duplicateId: 0
%% body:
可见光的频率范围从\tkanswer[2.6cm]{$3.9\times10^{14}$（或 $4\times10^{14}$）}赫兹到\tkanswer[2.6cm]{$7.9\times10^{14}$（或 $8\times10^{14}$）}赫兹。光合作用是地球上生物生存、繁荣和发展的根本源泉。高等植物叶绿体中的色素，主要吸收\tkanswer[1.2cm]{红橙}光和\tkanswer[1.2cm]{蓝紫}光。钠元素在一定条件下发出\tkanswer[0.8cm]{黄}色光，利用此性质，可制成公路上照明使用的高压灯。

\onepicture[w=3.2cm, align=right]{figs/45.png}

%% answer: 
%% memo:
\memoanswer{
本题原卷及材料中未提供详解，待补充。
}

\item
%% number: 50
%% paperName: 2003年上海理综第 50 题 6 分
%% typeId: 1
%% type: 单选题
%% chapter: 运动和力
%% point: 牛顿第一定律
%% method: 无
%% score: 6
%% degree: 650
%% duplicateId: 0
%% body:
理想实验有时更能深刻地反映自然规律。伽利略设想了一个理想实验，其中有一个是经验事实，其余是推论。

①减小第二个斜面的倾角，小球在这斜面上仍然要达到原来的高度

②两个对接的斜面，让静止的小球沿一个斜面滚下，小球将滚上另一个斜面

③如果没有摩擦，小球将上升到原来释放时的高度

④继续减小第二个斜面的倾角，最后使它成水平面，小球要沿水平面作持续的匀速运动

\onepicture[w=8cm, align=right]{figs/50.png}

请将上述理想实验的设想步骤按照正确的顺序排列\tkanswer[2.2cm]{②③①④}（只要写序号即可）。

在上述的设想步骤中，有的属于可靠的事实，有的则是理想化的推论。下列关于事实和推论的分类正确的是\xzanswer{B}

\fourchoices[answer=B]
{①是事实，②③④是推论}
{②是事实，①③④是推论}
{③是事实，①②④是推论}
{④是事实，①②③是推论}

%% answer: B
%% memo:
\memoanswer{
本题原卷及材料中未提供详解，待补充。
}

\end{enumerate}

\end{document}
